\section{Afterbody loads, body segmentation and the Iris sustainer}
\label{sec:afterbody}
A body tube can occupy the same space whether a designer draws it as one piece or splits it at a coupler. In the previous research blend, the nose chart read just the first tube length. On a 12-inch ogive with a 48-inch, 4-inch-diameter cylinder, splitting the cylinder into four pieces changed the predicted CP by up to $0.696D$ at zero incidence. Joining equal-diameter, touching tubes for the chart lookup reduced the maximum difference below $10^{-8}D$ over all 453 tested Mach/incidence combinations. This correction addresses the representation of the same exterior in the model.

\sub{A second hypothesis concerns the taper at the base.} A contracting boattail supports a negative normal-force increment. Even a small error in this load matters when it acts many diameters behind the nose. On a bare body the positive forebody and negative tail loads can nearly cancel. CP can then lie forward of the nose and be especially sensitive to errors in either load. For this reason we test the force and first moment of the taper against a straight control as well as the CP of the whole rocket.

\sub{Methods tested.} The prior blend fixes the fin and nose chart weights at 0.32 and 0.40. We evaluated three additional methods individually and jointly. The code extracts the original Fortran DATA arrays from the archived \href{https://www.pdas.com/datcomdownload.html}{USAF Stability and Control Digital DATCOM distribution}, subroutine \texttt{SUPBOD}. Every new table uses multilinear interpolation with the first Fortran index contiguous. No coefficient was fitted to the new rockets.

\begin{enumerate}[leftmargin=*]
\item \textbf{Center section.} For an eligible nose of length $L_n$ and base diameter $D_n$, group the continuous, same-stage center section up to the last eligible tube. The extended model permits monotone conical steps whose diameters remain within 10\% of $D_n$. With end diameter $D_1$, set $\bar D=(D_n+D_1)/2$, $u=\sqrt{M^2-1}\,D_n/L_n$, and $f=\min(5,L_cD_n/(L_n\bar D))$. Read normal slope and CP fraction from figures 4.2.1.1-21A/B and 4.2.2.1-23A/B. Normalize slope on the area of $D_1$ and multiply CP fraction by $L_n(1+f\bar D/D_n)$. Blend normal force and its nose-referenced first moment. Retain the unblended fraction of the incorporated step loads. Applying this smooth-body construction across short stepped transitions is an engineering approximation. The equal-tube variant stops at the first change in radius or non-tube component.
\item \textbf{Terminal taper.} For taper length $L$, fore/aft diameters $d_1,d_2$, $q=d_2/d_1$, $\beta=\sqrt{M^2-1}$ and taper half-angle $\theta$, use the fore-area coefficient. For $q<1$, figure 4.2.1.1-22A gives $\Delta N=(1-q^2)F(L/(\beta d_1))$. Figure 4.2.2.1-24 gives its station from the start of the taper, $x_t=L H(\beta\tan\theta,q)$. For a flare, figure 4.2.1.1-22B gives $\Delta N=G(q)\cos^2\theta$. The tabulated moment $K$ from figure 4.2.2.1-25A gives $x_t=-K(\theta,1/q)d_1q^3/\Delta N$. Convert area, shift the moment to the rocket datum and blend force and moment. The boattail port uses the chart's local load station and translates it by the taper start in OpenRocket. This is a component implementation of the charts.
\item \textbf{Crossflow.} Replace the constant 1.1 in the symmetric body's finite-incidence lift term by the crossflow coefficient in figure 4.2.1.2-35B evaluated at $M|\sin\alpha|$. Keep OpenRocket's projected area, centroid and $\sin^2\alpha$ term. This correction vanishes at zero incidence.
\end{enumerate}

\sub{Eligibility.} Apply the established smooth envelope over Mach 1.1--1.5 and 4.65--5.0, with full incidence weight through $3^\circ$ and zero at $8^\circ$. The taper must be the last active centerline part and touch a radius-matched tube. Use chart limits $\beta\tan\theta\leq0.5$ and $L/(\beta d_1)\leq2.5$ on contractions. Smoothstep to zero over the final tenth of each interval. An abrupt switch at the first limit was visible on the ARCAS Mach sweep before this correction. Flare queries require $1<q\leq3$ and $5^\circ\leq\theta\leq90^\circ$. The stage mask invalidates cached component eligibility when a booster separates.

\clearpage
\sub{Native RASAero campaign.} RASAero II 1.0.2.0 in the Windows 11 VM produced independent Aero Plots exports for 13 configurations. Three existing native Iris component exports complete the comparison. Each row below uses 151 Mach values from 1.50 to 3.00 in 0.01 increments at each incidence. The reference is the native angle-specific CP in inches aft of the nose, normalized by that configuration's maximum body diameter. P denotes the prior blend. N joins equal tubes and applies DATCOM only to terminal contractions. A includes the stepped-center, terminal taper and crossflow methods together.
\begin{center}\small\setlength{\tabcolsep}{4.5pt}\begin{tabular}{l rrr rrr rrr}\toprule
&\multicolumn{3}{c}{$0^\circ$}&\multicolumn{3}{c}{$2^\circ$}&\multicolumn{3}{c}{$4^\circ$}\\
Case & P & N & A & P & N & A & P & N & A\\\midrule
\tabinput afterbody-static.tex \bottomrule\end{tabular}\end{center}

\sub{Geometry controls.} A has a 12-inch tangent ogive and 48-inch afterbody of nominal diameter 4 inches. Its 4-inch terminal segments end at 3.0 or 4.8 inches. Its internal-step alternative has a 36-inch tube, 0.50-inch step to 4.2 inches, then an 11.5-inch tube. B uses a 9-inch conical nose and 36-inch, 3-inch-diameter afterbody. Its 3-inch tapers end at 2.4 or 3.6 inches. The three-fin B pair shares the same 4.5/2-inch root/tip, 2.25-inch span, 2.5-inch sweep, 0.125-inch square section and fore-aft placement. Family B was reserved before the A/B calculations. After inspecting those results we fixed the narrower method and requested a fresh confirmation pair C. C has a 5-inch diameter, 18.5-inch tangent ogive, and 49-inch afterbody. In the variant the final 6 inches taper to 4.5 inches. Four fins of 6/2-inch root/tip chord, 3.5-inch span, 4.5-inch sweep and 0.15-inch thickness occupy the same root stations, 54--60 inches, on both models.

On the straight A and B controls, the crossflow coefficient is the only active change in the all-method model. For A it improves CP RMS from 0.33 to $0.24D$ at $2^\circ$ and worsens it from 0.38 to $0.50D$ at $4^\circ$. B shows the same reversal. This incidence dependence argues for retaining the original crossflow law in the narrow variant.

The straight Iris controls remove the steps while retaining the 136.5-inch length, 4.02-inch cylinder and original fin root stations at 128.5--136.5 inches. The static reconstruction uses the preceding component's aft radius as the front of a native transition and \texttt{Transition/Diameter} as its outlet. In two additional VM exports, changing only \texttt{RearDiameter} left the coefficients identical. Changing only \texttt{Diameter} moved the rear OD in the editor and altered all 900 checked CP values. A separate \texttt{BoatTail} stores its outlet in \texttt{RearDiameter}. Equal values in the two serialized diameter fields can describe a real taper. This resolves an ambiguity in the earlier literal-CDX1 diagnostic. Each terminal \texttt{BoatTail} contributes geometry once. The body-tube boattail fields duplicate its dimensions. The OpenRocket static fixtures reproduce the external axisymmetric profile and attached fins. The preceding native camera-area removal controls gave identical normal loads and CP at all 906 matched states. The integrated flight uses the actual ORK with its pods and motors.

\clearpage
\sub{Force and moment of a geometry change.} For each matched pair use the same $D_*$ and area $S_*$ for both configurations. At finite incidence recover OpenRocket $C_N=\alpha N$ with $\alpha$ in radians. For each code separately compute
\[
 \Delta C_N^*=\left[C_N(D/D_*)^2\right]_{\rm variant-control},\qquad
 \Delta J^*=\left[C_N(D/D_*)^2x_{\rm CP}/D_*\right]_{\rm variant-control}.
\]
The table gives RMS error in these increments at $2^\circ$. Use $D_*=4$, 3, 5 and 4.20 inches for A, B, C and Iris. These nose-referenced first moments use aft-positive station. The same analysis at $4^\circ$ and all component-only ablations are in the result CSVs.
\begin{center}\small\begin{tabular}{lrrr rrr}\toprule
&\multicolumn{3}{c}{RMS error in $\Delta C_N^*$}&\multicolumn{3}{c}{RMS error in $\Delta J^*$}\\
Variant & P & N & A & P & N & A\\\midrule
\tabinput afterbody-pairs.tex \bottomrule\end{tabular}\end{center}

\sub{Wind-tunnel evidence.} ARCAS contains a contracting afterbody. Across its five supersonic measurements the candidate changes CP RMS from 0.312 to $0.363D$, normal-force-slope RMS from 2.63 to 2.31/rad, and pitching-moment-slope RMS from 4.94 to 3.08/rad about the benchmark datum. Yuska B has a straight afterbody and gives the same $0.136D$ CP RMS. These measurements give a mixed assessment of boattail transfer.

For a separate flare check we used Jorgensen, Spahr and Hill's \href{https://ntrs.nasa.gov/citations/19720065133}{NASA TM X-653}, figures 1(b) and 4(a,b). Both configurations have a 1.25-inch-diameter midbody, blunted $14^\circ02'$ cone, a $2D$ cylinder and a $1.37D$ terminal segment. One terminal segment is straight and the other flares to $1.48D$. The numerical forebody uses a sphere of radius $0.25D$ tangent to the cone, with 40 frusta over the spherical portion. We read the plotted small-angle slopes and CP at approximately Mach 1.97, 2.94 and 4.06. Normalize both force slopes to the midbody area, convert degrees to radians, and measure CP and $J$ from the nose--cylinder junction. The flare's total CP RMS improves from about 0.074 to $0.053D$. The measured flare-minus-straight increment gives the opposite assessment. Its slope RMS rises from 1.06 to 1.10/rad and its first-moment RMS rises from 2.26 to 2.57/rad. The straight control already has $0.95D$ CP RMS error. The original plot and digitization are archived. Two decimal places describe graph readings only.

\sub{Checks on the implementation.} Each fixed-incidence case uses a single solver across its Mach grid. A detach/reattach test compares a reused calculator with a fresh one at each stage state. The earlier ARCAS and Yuska B fixtures were scanned every 0.001 Mach from 0.750 to 5.001, checking finite loads and adjacent changes in CP and slope. The packaged narrow JAR reproduces 8,613 harness rows with a maximum CP difference below $2\times10^{-12}$ inches compared with the class-directory run. Legacy flight histories for both research runs are byte-identical to the previously archived legacy flight. The prior-mode Iris replay reproduces the earlier shared-state solver output to $10^{-9}$. The segmented cylinder's one-, two- and four-tube representations also agree in normal slope to $10^{-8}$/rad with the geometry correction.

\clearpage
\sub{Iris in two complete flight solvers.} The preferred first transition is 0.50 inches long. We reran the full two-stage flight in the Windows VM with the original O6836 and N3366 motors. All 16,664 rows and all 24 exported columns reproduced the archived trajectory byte for byte. The OpenRocket results below are separately integrated flights using identical seeds. Within the 15.52--20.77 s sustainer burn we interpolated the separated-stage OpenRocket histories at the 526 native timestamps. Minimum and its time use the integration samples including interpolated endpoints. RMS and bias are in calibers and take OpenRocket minus RASAero.
\begin{center}\small\begin{tabular}{lrrrrr}\toprule
Model & Burnout & Minimum & Time of min (s) & Burn RMS & Burn bias\\\midrule
\tabinput afterbody-flight.tex \bottomrule\end{tabular}\end{center}
\fig{.98}{afterbody-comparison}{Left, zero-incidence CP error against native Aero Plots. Right, integrated Iris stability through the sustainer burn. The narrower variant overlaps the prior blend to plotting precision on this flight.}

\sub{Attribution on Iris.} Replaying the same native Mach, incidence and stage at burnout reduces the CP error from $+0.572D$ to $+0.112D$. The stepped-center method alone gives $+0.112D$, the taper alone $+0.572D$ and the crossflow alone $+0.572D$. The complete sustainer ends in a tube. On the 975 sampled sustainer states above Mach 2 the all-method CP RMS drops from $0.680D$ to $0.285D$. Yet on the stepped finless sustainer the zero-incidence static CP RMS rises from $2.86D$ to $4.16D$. At $2^\circ$ the matched Iris step's first-moment discrepancy also grows. The lower assembled burnout margin therefore cannot establish the assumed pressure distribution along the steps. The minimum remains about 2.6 s early and is now 0.23 cal below RASAero.

\sub{Use in MIT 6.6.} Aggregate geometrically identical tubes. Offer contracting boattails as an experimental comparison for configurations within the chart envelope, with the ARCAS CP regression visible. Keep the stepped-center method available to bound model sensitivity on Iris. Keep the previous flare and crossflow terms for general use. The separate patch and JAR in \texttt{implementation/} select the narrow method by default when the experimental fin option is active. In a fresh JVM use \texttt{-Dstudy.afterbody=} followed by \texttt{previous}, \texttt{segments}, \texttt{center}, \texttt{tail}, \texttt{crossflow}, \texttt{combined}, or \texttt{narrow}. All modes preserve the published fin-blend constants. The tables are component ports of selected DATCOM methods, with a geometry extension expressly identified above. The comparisons cover static longitudinal loading at small incidence.
